% 【专题研讨】所属：第2章（由 chapters/revised/02-skepticism-luck.tex \input 引入）
\minitopic{专题研讨：认识风险、能力与鲁棒性}\index{认识风险}\index{过程个体化} 反运气理论可被理解为认识风险理论：主体在当前证据和方法下接受$p$，承担多大出错风险？风险不只是一项概率。错误发生的接近性、损失、错误是否可发现、是否由主体能力控制，都影响评价。高概率真信念可能仍在个案上缺乏合适联系。 概率风险关注频率或置信，模态风险关注近邻世界，能力风险关注成功来源，责任风险关注主体是否响应警报。一个系统可以统计准确却在关键子群不安全，也可以个案安全却因使用者毫无理解而缺少成就。综合分析应保留四栏。

\minitopic{敏感性与安全性的系统比较}敏感性条件检验$p$为假时主体是否仍相信$p$；安全性检验主体以相同方式相信$p$时$p$是否仍真。两者的起点不同。敏感性强调随事实变化，安全性强调排除近邻错误。把它们都叫“可靠”会隐藏彩票和闭合差异。 彩票在开奖前是典型压力：若这张票中奖，主体仍基于极低概率相信会输，所以不敏感；但在许多近邻世界票确实输，信念可能显得安全。若安全区域包含中奖世界，答案又变化。案例显示接近性不能只由频率决定。 必要真理对敏感性构成空真问题。若$p$不可能为假，就没有反事实测试；随便猜到数学真理似乎也敏感。安全性同样可能因真理必然性自动满足。两种条件都需与适当方法或能力结合，不能单独定义所有知识。 闭合压力方向相反。主体敏感地知道“这是斑马”，推知“不是伪装骡子”，在伪装世界仍会相信后者，故结论不敏感。安全性更易保留某些闭合，但复杂推理可能改变形成方式。任何“保留闭合”的结论都需限定主体胜任推理且无新击败者。

\begin{center}
\begin{tabularx}{\textwidth}{@{}p{2.3cm}XX@{}}
\toprule
检验 & 固定内容 & 主要压力\\
\midrule
敏感性 & 让$p$为假，观察信念是否改变 & 必然真理、彩票、闭合失败\\
安全性 & 保持相关信念方法，观察$p$是否仍真 & 近邻范围、过程个体化、低概率故障\\
能力归因 & 保持成功，改变能力贡献 & 环境运气、工具与团队功劳\\
\bottomrule
\end{tabularx}
\end{center}

\minitopic{过程个体化与事前标准}同一事件可描述为“使用视觉”“在雾中看灯”“周二读一块红屏幕”“在这次真命题出现时阅读”。越具体越容易获得完美可靠率，越宽泛越可能混入无关失败。过程类型若在看到结果后选择，理论可以让任何信念可靠或不可靠。 事前个体化要求类型对应可重复机制、主体功能和错误模式。视觉与记忆有不同心理机制；同一检测的两种试剂有不同故障；同一网页在正规数据库和伪造镜像中环境不同。类型应能预测新案例，而非把真值写入名称。 任务也限制类型。辨认常见物体、识别罕见肿瘤和在假立面环境辨认谷仓，不是同一难度领域。能力评价需要说明对象范围和正常条件。把“正常”留空会把所有反例排除为不正常；应由训练、设计和实际使用规范给出。 群体和工具使个体化更复杂。医生读影像的过程包括设备校准、算法增强、视觉训练和会诊。只评价医生眼睛会漏掉上游故障，评价整个医疗系统又可能太宽。可以按问题分层：个体信念形成、团队诊断流程、机构发布分别有过程类型。

\minitopic{自我验证与虚假安全}一个不可靠仪器若只用自己的输出校准，可能每次都与自身一致。主体据此相信它可靠，再把可靠性信念用于证明读数知识，形成自举。内部一致提高，却没有独立世界约束。 自举说明“获得许多真读数”不足以自动产生来源可靠性知识，尤其当主体无法独立知道哪些读数真。可加入外部标准、异质仪器、盲样本和故障注入。每项都让来源在事实变化时受到新约束。 但完全独立验证也不总可能。新仪器测量前所未见对象，没有同类金标准；科学家依赖理论、部件测试和间接预测组成证据网。合理要求是减少共同错误、让关键假设分散，而非要求无任何共享前提。 虚假安全还来自过窄测试。模型在与训练数据近似的测试集稳定，却在新医院失败；团队把“测试附近不易错”误当真实部署安全。模态区域必须包含任务实际可能变化，而不只包含开发者方便模拟的变化。

\minitopic{德性认识论中的能力生态}能力不是主体内部孤立模块。阅读能力依赖光线和文本规范，专家判断依赖训练与反馈，证言判断依赖社会记录。生态观点把能力看成主体策略与环境信息结构的匹配。环境改变后，同一习惯可能从德性变成脆弱启发式。 德性可以是最低知识条件，也可以是更高探究理想。正常视觉产生普通知识，不要求主体具备非凡认真；开放、勇气和谦逊则改善长期探究。将二者混合，会因一个人性格有缺陷而否定其所有知觉知识。 能力归功有程度。射手命中可能七成归技术、三成受风；科研发现由个人洞见、仪器和同事共同促成。知识理论不必给精确百分比，但要判断偶然因素是否主导。若成功解释主要诉诸错误抵消，就不是恰当成就。 反思德性包括知道何时不使用一阶能力。优秀医生不仅识别常见病，也知道病例何时超出训练；优秀信息使用者不仅快速判断，也知道来源何时需要核验。边界识别是能力的一部分，不是缺乏自信。

\minitopic{认识风险的制度分配}高风险系统常把最终决定交给个人，却把信息、时间和工具控制留在组织。若个人只能在数秒内批准模型建议，说“最终由人负责”没有创造能力或安全。制度应把风险分配到能够改变原因的角色。 开发者负责数据和测试范围，部署者负责环境匹配，管理者负责资源与升级，使用者负责个案警报。事故发生后应沿因果和权限追踪，而非只惩罚最后点击者。责任结构本身影响信息是否上报。 冗余要异质。同一算法运行三次不会发现系统偏差；不同数据、机制和团队更能降低共同失败。异质冗余成本高，可集中于损失大、不可逆和模型外样本。认识论在此帮助解释哪些错误世界应被设计纳入。 风险沟通也构成安全。系统给出置信区间、适用范围和已知失效，使用者才能适当响应；只显示一个绿色结果会隐藏击败者。透明不是提供全部源代码，而是提供当前判断所需的可行动信息。

\minitopic{比较视角：训练、环境与解蔽}《荀子》关于学习、师法和“解蔽”的讨论把认识能力放在长期训练、语言和公共校正中\parencite{xunzi2014}。这与生态德性观有问题层面的相似：个体不会仅凭内在天赋自然形成稳定判断，环境和实践塑造什么被看见。 但古典“蔽”包含欲望、立场与政治秩序，不能缩成统计错误率。现代安全性可以精确描述近邻失败，却可能忽略主体为何固执选择某一信息环境；修养论强调改变注意和欲望，却未提供形式化风险度量。两者互补在问题上，不等同于同一理论。 师法也有双刃。可靠教师与规范可降低初学者风险，权威僵化则让群体共同复制错误。能力生态必须包含权威可纠正、异议可表达和学习者逐步获得独立判断，而非以传统稳定性代替可靠性。

\minitopic{综合研讨：自动驾驶接管}自动驾驶在高速路通常比疲劳司机可靠，但遇到施工反光锥时置信下降，要求司机在五秒内接管。司机长时间未操作，界面只发一个提示音。事故前系统正确识别风险，组织却没有提供真正可用的人类能力。 若只看平均准确率，系统可靠；若看该施工分布，错误率高；若看接管信念，司机在近邻情形容易误判；若看能力归因，五秒监督不是稳定技能。四种评价揭示不同风险。 改进可包括更早预警、持续驾驶员状态监控、在无法安全接管时自动降速、施工场景专项测试、保留人机分歧日志。每项改变不同近邻世界和责任节点。简单增加“请注意”文字未必改善能力。 若事故没有发生，系统恰好驶过施工区，也不能仅以成功结果证明安全。反事实审计要问小幅变化会怎样，成功是否由设计机制保证，司机是否有真实控制。反运气认识论在这里转化为工程鲁棒性和制度问责。

\begin{tcolorbox}[breakable,colback=white,colframe=blue!30!black,title={专题研讨任务},fonttitle=\bfseries]
选择一个人机或科学系统，分别评估频率风险、敏感性、安全性、能力归因和责任风险。定义事前过程类型，设计一次故障注入与异质复核，再说明哪项成功仍可能只是认识运气。
\end{tcolorbox}
